\documentclass[11pt,a4paper]{article}
\usepackage[margin=1.7cm]{geometry}
\usepackage{fontspec}
\usepackage{enumitem}
\usepackage{titlesec}
\usepackage[hidelinks]{hyperref}
\usepackage{xcolor}
\usepackage{tabularx}
\usepackage{setspace}
\setmainfont{PingFang SC}
\linespread{1.12}
\pagestyle{empty}
\setlength{\parindent}{0pt}
\setlist[itemize]{leftmargin=1.6em,itemsep=1pt,topsep=1pt,parsep=0pt,partopsep=0pt}
\setlength{\emergencystretch}{2em}
\titleformat{\section}{\large\bfseries}{}{0em}{}[\vspace{-0.4em}\titlerule]
\titlespacing*{\section}{0pt}{0.8em}{0.5em}
\newcommand{\entry}[4]{%
  \textbf{#1} \hfill #2\\
  \textit{#3} \hfill \textit{#4}\\
}
\begin{document}

{\LARGE\textbf{田昊琦}}\\[4pt]
\href{tel:19357571423}{19357571423} \;|\; \href{mailto:cathytian@zju.edu.cn}{cathytian@zju.edu.cn} \;|\; 求职意向：产品经理

\section*{教育背景}
\entry{浙江大学}{2025.09 -- 2027.06}{金融（数字金融方向）研究生}{课程均分90+}
\entry{浙江大学}{2021.09 -- 2025.06}{金融学本科}{排名前30\%，保送研究生}

\section*{实习经历}
\entry{京东}{2026.05 -- 至今}{AI 数据产品经理}{}
\begin{itemize}
  \item AI 平台能力建设：参与 Data+AI 平台（对标 BI 产品）建设，负责提示词体系与记忆体系设计，支撑 AI 问数到报告生成全链路。
  \item 数据体系与工具化：设计自动驾驶业务数据指标体系与埋点规范；搭建自动化工具（基于 PRD 与原型图自动生成埋点方案，并自动完成埋点录入与指标口径填写），降低重复性数据工作成本。
  \item AI 智能分析闭环：搭建自然语言意图提取、API 取数、报告生成、在线文档填入的全链路 Agent，大幅压缩取数分析时间。
\end{itemize}

\entry{滴滴出行}{2025.08 -- 2026.03}{数据中台产品经理实习生}{}
\begin{itemize}
  \item 标签逻辑重构：负责地理数据产品核心逻辑迭代，设计标签有效期与互斥逻辑，将非标需求转化为通用配置能力，提升运营配置效率 50\%，每周节省工时约 2 小时。
  \item 清淤决策优化：迭代清淤需求，通过热力图可视化与标签自动上下标，将维护频率从每周 2 次提升至每日实时，提升数据鲜活性，每周释放运营工时约 40 分钟。
  \item 系统治理与运营提效：主导批量撤回功能设计，解决跨系统数据一致性问题，累计处理点位 4000+ 个，日均节省运营人力约 40 分钟（峰值期日省 9 小时）。
  \item 构建围栏运营指标体系：围绕系统效率、操作行为与点位状态三大维度设计 15+ 核心指标，支撑自动化建点闭环；提出采纳率漏斗模型量化转化效果并识别关键瓶颈；搭建多维可视化看板支持城市与点位类型下钻。
  \item 平台智能化：基于 Dify 搭建标签相似度识别、标签命名、标签审核工作流，并设计平台融入方案。
\end{itemize}

\entry{阿里巴巴集团-钉钉}{2025.03 -- 2025.06}{AI 表格产品实习生}{}
\begin{itemize}
  \item 产品功能迭代：深入 ToB SaaS 业务场景，总结 VOC 挖掘功能需求；参与“主键即文档、查询页面、AI 生图”等功能 PRD 撰写、评审与验收。
  \item Agent 搭建：主导产品文档知识库撰写，通过意图识别与工作流编排，配置具备需求收集、智能问答、自主学习能力的 AI 助理。
\end{itemize}

\section*{项目经历}
\entry{LLM 大模型可视化}{2026.03 -- 至今}{独立开发者}{}
\begin{itemize}
  \item 产品 Demo：\url{https://visualize-llm.netlify.app/}
  \item 需求洞察与产品定义：针对模型使用者“黑盒认知”痛点，从 0 到 1 规划基于 WebGPU 的可视化工具。
  \item 核心架构与交互设计：规划 Token 实时渲染矩阵、Top-K 概率图、UMAP 语义降维空间等模块；创新“逐步干预模式”解释“模型幻觉”成因，并借助 Claude Code 快速验证高保真交互原型。
\end{itemize}

\entry{浙江大学 × 蚂蚁集团产品经理训练营}{2024.11 -- 2024.12}{产品负责人}{}
\begin{itemize}
  \item 0-1 产品定义与团队协同：作为产品创意发起人及队长，主导从用户痛点洞察到功能方案推导，组织跨专业团队协作并推动 PRD、故事线与关键节点对齐，确保按期高质量交付。
  \item 高保真原型与成果呈现：独立完成 Figma 高保真原型与核心交互流程设计，并主导演示路演，清晰传递产品价值、用户体验与商业逻辑，获训练营优秀个人。
\end{itemize}

\section*{技能与爱好}
\begin{itemize}
  \item 软件技能：Claude Code、Coze、Dify、Microsoft Office、Photoshop、Canva、SQL、Python、Figma、Trae。
  \item 语言能力：汉语普通话（母语），英语（六级 550+）。
  \item 爱好特长：志愿活动（180+ 小时）、脱口秀（多次表演经历）、小说写作（豆瓣签约作品 1 本）。
\end{itemize}

\end{document}
